\documentclass{article}
\usepackage{graphicx} % Required for inserting images
\usepackage[margin=2cm, top=3cm]{geometry} % margine superiore ridotto a 1 cm
\usepackage{listings} % Libreria per Python
\title{Intelligenza Artificiale e Machine Learning}
\date{Ottobre 2024}
\begin{document}

\maketitle

\section{Introduzione al linguaggio Python}
\subsection{Le basi}
print() per stampare\\
           type() per conoscere il tipo dell'argomento\\
           Non serve specificare il tipo della variabile\\
           Per avere solo la parte intera della divisione //\\
           Per inserire dati da input basta chiamare la funzione input()\\
           In python esiste una modalità interattiva, siamo in modalità interattiva se prima ci sono . . .\\
           E possibile gestire le eccezioni in questo modo:
           \begin{lstlisting}[language=Python, frame=single, numbers=left, caption=Esempio di eccezioni in codice Python]
try:
prompt= 'Dammi un dato!'
dato=input(prompt)
print(dato)
int(dato) per effettuare una conversione
except:
print('Please enter a number')
\end{lstlisting}

\subsection{Le condizioni}
Gli if hanno questa sintassi:
\begin{lstlisting}[language=Python, frame=single, numbers=left, caption=Esempio di condizioni in codice Python]
     if x>0: 
          print('x e' positivo')
          else:
          print('x e' negativo')
\end{lstlisting}
\subsection{Funzioni e funzioni built-in}
Esistono funzioni built-in come max() e min()\\
Esiste il modulo math 'import math' per eseguire la maggior parte delle funzioni matematiche es. math.sin(x), le funzioni trigonometriche accettano un argomento in radianti\\
Quindi per convertire gradi/360 * 2pi\\
Il modulo random permette di avere funzioni che generano numeri pseudo-casuali\\
è possibile definire funzioni con def funzione(): istruzioni\\
Non si possono creare variabili con lo stesso nome di una funzione dato che in python la funzione è un oggetto funzione\\
Se si tenta di assegnare a una variabile il valore di una fruitful function si otterrà un valore speciale chiamato "None"\\
\subsection{I cicli}
Qui non si usano parentesi graffe, infatti per capire il blocco di un istruzione condizionale, una funzione o un ciclo si identa semplicemente\\
 Nel codice, quando sono finite le istruzioni si toglie l'identamento\\
 Si può fermare il ciclo quando si vuole con l'istruzione break\\
 Contrariamente si può usare l'istruzione continue per saltare alla successiva iterazione senza terminare il corpo del ciclo per l'iterazione corrente\\
 Esiste il for each:\\
 \begin{lstlisting}[language=Python, frame=single, numbers=left, caption=Esempio di for-each in codice Python]
lista = [1, 2, 3]
for numero in lista:
    print(numero)
print('Done')
\end{lstlisting}
Per trovare il valore più grande in una lista:
 \begin{lstlisting}[language=Python, frame=single, numbers=left, caption=Esempio di for-each in codice Python]
largest = None
for itervar in [3, 41, 12, 9, 74, 15]
if largest is None or itervar >largest:
    largest = itervar
print(largest)
\end{lstlisting}
Ovviamente il codice per il più piccolo è pressochè identico
\subsection{Le stringhe}
Funzione len per la lunghezza di una stringa\\
è possibile ottenere delle sottostringhe tramite i :
 \begin{lstlisting}[language=Python, frame=single, numbers=left, caption=Esempio di operatore slice in codice Python]
s = 'Monty Python'
print(s[0:5]) # nota: l'ultimo numero (in questo caso 5) e' escluso
>>Monty 
\end{lstlisting}
\noindent
Se si omette il primo indice, la sottostringa parte dall'inizio della stringa, se si omette il secondo indice, la sottostringa arriva alla fine della stringa
\end{document}
